\unitopen{27}
\begin{candos}
\item użyć dopełniacza (\de{Genitiv}) przy rzeczowniku i po przyimkach \de{wegen, trotz, statt, während};
\item dobrać końcówkę przymiotnika bez rodzajnika (\de{frischer Kaffee, gute Nachrichten});
\item zrozumieć nagłówek prasowy i zdanie z dopełniaczem;
\item zapytać \de{Wessen \dots?} i odpowiedzieć.
\end{candos}

\section{Słownictwo}
\vocabtable{27}

\section{Dopełniacz}
Dopełniacz oznacza posiadanie i występuje po kilku przyimkach. Rodzajniki: \artm{des} (m, n), \artf{der} (f, mnoga). Rzeczowniki rodzaju
męskiego i nijakiego dostają \de{-s} lub \de{-es}: \de{das Ende des Jahres}, \de{der Name des Kunden}, \de{die Meinung des Chefs}.

\begin{center}\small
\begin{tabular}{@{}lll@{}}
\toprule
\textsf{\textbf{Rodzaj}} & \textsf{\textbf{Przykład}} & \textsf{\textbf{Uwaga}}\\
\midrule
m & \de{das Auto des Mannes} & \de{-es} po jednosylabowych, \de{-s} po dłuższych\\
f & \de{die Idee der Kollegin} & bez końcówki rzeczownika\\
n & \de{der Preis des Produkts} & \de{-s}\\
mnoga & \de{die Meinung der Kollegen} & bez końcówki\\
\bottomrule
\end{tabular}
\end{center}
Pytanie: \de{Wessen Laptop ist das?} (czyj?). W mowie dopełniacz zastępuje się często \de{von} + celownik:
\de{das Auto von meinem Bruder}. Przyimki z dopełniaczem: \de{wegen} (z powodu), \de{trotz} (mimo), \de{statt} (zamiast),
\de{während} (podczas), \de{innerhalb} (w obrębie), \de{außerhalb} (poza): \de{wegen des Wetters}, \de{trotz der Verspätung}.

\section{Przymiotnik bez rodzajnika}
Bez rodzajnika przymiotnik niesie końcówkę rodzajnika: \de{frischer Kaffee} (m), \de{frische Milch} (f), \de{frisches Brot} (n),
\de{frische Brötchen} (mnoga). Celownik: \de{mit frischem Kaffee}, \de{mit frischer Milch}. Mnoga w mianowniku:
\de{gute Nachrichten}, \de{neue Kollegen}.

\begin{falle}
Po polsku dopełniacz stoi \emph{po} rzeczowniku („rodzina brata”), po niemiecku też: \de{die Familie des Bruders},
nie \de{die Familie dem Bruder}. Błąd: pominięcie \de{-s} (\de{das Auto des Mann}). Drugi: \de{wegen} z celownikiem
(\de{wegen dem Wetter}) jest potoczne, ale w piśmie błędne. Trzeci: pamiętaj o wyjątkach: \de{der Name $\rightarrow$ des Namens}
(rzeczowniki na \de{-n}), \de{der Herr $\rightarrow$ des Herrn}.
\end{falle}

\begin{szwajcaria}
W Szwajcarii dopełniacz jest częstszy w piśmie urzędowym (\de{des Kantons}, \de{des Bundes}) niż w mowie. Nagłówki prasowe są
w standardzie (\de{NZZ}, \de{Tages-Anzeiger}, \de{SRF}); do ćwiczenia czytania wybierz \de{SRF News}, bo język jest przystępny.
\end{szwajcaria}

\begin{praca}
\de{Wegen des Ausfalls des Servers} \de{konnten wir nicht arbeiten.} \de{Trotz der Verspätung ist das Projekt pünktlich fertig.}
\de{Die Anforderungen des Kunden haben sich geändert.} \de{Statt eines Meetings schreiben wir eine kurze Zusammenfassung.}
\end{praca}

\section{Dialogi}
\textbf{Dialog 1 --- omawianie wiadomości}
\begin{dialog}
\turn{A}{Hast du die Nachrichten von heute gelesen?}{Czytałeś dzisiejsze wiadomości?}
\turn{B}{Ja, wegen der hohen Energiepreise gibt es eine neue Diskussion.}{Tak, z powodu wysokich cen energii trwa nowa dyskusja.}
\turn{A}{Und was meint die Regierung?}{A co sądzi rząd?}
\turn{B}{Die Regierung sagt, dass sie trotz der Krise keine neuen Schulden machen will.}{Rząd mówi, że mimo kryzysu nie chce zaciągać nowych długów.}
\end{dialog}

\section{Ćwiczenia}
\begin{cwiczenia}
\exitem{Dopełniacz: \de{das Ende (das Jahr)}.}{\de{das Ende des Jahres}}
\exitem{Dopełniacz: \de{die Meinung (der Chef)}.}{\de{die Meinung des Chefs}}
\exitem{Dopełniacz: \de{die Idee (die Kollegin)}.}{\de{die Idee der Kollegin}}
\exitem{Dopełniacz mnogi: \de{die Meinung (die Kollegen)}.}{\de{die Meinung der Kollegen}}
\exitem{Przyimek: \de{\dots{} des Wetters bleibe ich zu Hause.} (z powodu)}{\de{Wegen}}
\exitem{Przyimek: \de{\dots{} der Verspätung kommen wir pünktlich.} (mimo)}{\de{Trotz}}
\exitem{Zapytaj: czyj jest ten laptop?}{\de{Wessen Laptop ist das?}}
\exitem{Przymiotnik bez rodzajnika: \de{frisch\dots{} Brot} (nijaki).}{\de{frisches Brot}}
\exitem{Przymiotnik bez rodzajnika: \de{gut\dots{} Nachrichten}.}{\de{gute Nachrichten}}
\exitem{Przymiotnik z celownikiem: \de{mit frisch\dots{} Milch}.}{\de{mit frischer Milch}}
\exitem{Popraw: \de{Das Auto des Mann.}}{\de{Das Auto des Mannes.}}
\exitem{Przetłumacz: „Mimo awarii serwera pracowaliśmy dalej.”}{\de{Trotz des Serverausfalls haben wir weitergearbeitet.}}
\end{cwiczenia}

\section{Trening słownictwa (generowany)}
\drill{27}

\weekend{Przeczytaj jeden krótki artykuł na \de{SRF News} lub \de{Deutsche Welle} i wypisz trzy zdania z dopełniaczem
i trzy z przymiotnikiem bez rodzajnika. Streść artykuł w pięciu zdaniach na głos.}
